% =====================================================================
%  R. Nielsen, HOW ONE FORMS A WORLD-VIEW
%  In: For Idea and Reality.  A Journal, edited by R. Nielsen and Rud. Schmidt.  Year 1873 (Bind 8).
%  Copenhagen: Otto B. Wroblewsky, Publisher; Printed by G. S. Wibe.  Bind 8, pp. 1-18.  (Hvorledes man danner sig en Verdensanskuelse, in For Ide og Virkelighed.)
%  ENGLISH TRANSLATION, from transcription.tex (Danish), the text of record.
%  GENERATED by ebook-method/for-idee-1870-73/tr_build.py --write from this template and the translators' parts in
%  ebook-method/for-idee-1870-73/tr/ (TRANSLATOR-BRIEF.md; terms: GLOSSARY.md, its key terms approved by Hans
%  2026-10-13/14).  Conventions: TRANSLATION-PLAYBOOK.md.  Page numbers of the print in the margin.
% =====================================================================
\documentclass[12pt,a4paper,openany]{book}
\usepackage[utf8]{inputenc}
\usepackage[T1]{fontenc}
\usepackage{libertinus}
\usepackage{libertinust1math}
\usepackage{textalpha}
\usepackage[english]{babel}
\usepackage[protrusion=true,expansion=true]{microtype}
\usepackage{setspace}
\usepackage[margin=1.3in]{geometry}
\usepackage{fancyhdr}
\usepackage{hyperref}
\hypersetup{pdftitle={How One Forms a World-View}, pdfauthor={Rasmus Nielsen}, hidelinks}
\setlength{\headheight}{15pt}
\pagestyle{fancy}
\fancyhf{}
\fancyhead[LE,RO]{\thepage}
\fancyhead[C]{\textit{R. Nielsen: How One Forms a World-View (1873)}}
\setstretch{1.3}
\usepackage{marginnote}
\renewcommand{\thefootnote}{\arabic{footnote})}   % the print numbers each page's notes 1), 2)
\newcommand{\opage}[1]{\marginnote{\footnotesize\textit{[#1]}}}
\newcommand{\essayhead}[2]{\begin{center}{\large #1}\ifx&#2&\else\\[0.3em]{\normalsize #2}\fi\\[0.5em]\rule{3em}{0.4pt}\end{center}\medskip\noindent}
\newcommand{\signature}[1]{\par\smallskip\hfill #1\hspace*{2em}\par}

\begin{document}

% the title page of Bind 8 (1873-for-idee-og-virkelighed.pdf), in English
\begin{titlepage}
\begin{center}
\vspace*{3em}
{\Huge For Idea and Reality.}\\[1em]
\rule{8em}{0.4pt}\\[1em]
{\large A Journal,}\\[0.8em]
{\small edited}\\[0.4em]
{\footnotesize by}\\[0.6em]
{\large R. Nielsen and Rud. Schmidt.}\\[0.4em]
\rule{8em}{0.4pt}\\[0.6em]
{\large Year 1873.}\\[0.6em]
\rule{8em}{0.4pt}\\[4em]
Copenhagen.\\[0.4em]
Otto B. Wroblewsky, Publisher.\\[1.5em]
{\small Printed by G. S. Wibe.}\\[3em]
{\small Translated from the Danish}
\end{center}
\end{titlepage}

\essayhead{How One Forms a World-View.}{}
% --- p. 1 ---
\setcounter{footnote}{0}%
\opage{1}What is a world-view? Something that forces itself
upon one, something that can be had for a good bargain, and something
that must be bought at the dearest price, according to how one takes
it. All people have a world-view, each his
own; they can resemble one another so strikingly that one
would swear they were woven in one pattern, but they
can also be highly different. Individuals can
disagree about many things and yet have the same world-view;
they can also agree about many things, and yet
disagree in world-view. It is important for society
to know how much agreement in this respect
is necessary if civil order is to be maintained;
for though every human being must indeed have the right to
see the world and life with his own eyes, yet if the disagreement
between individuals goes to extremes, society
is in danger. There are real world-views that
cannot be reconciled; the thief, for example, has one world-view,
the police have another; two such cannot be ``mediated.''

Were it not for the sake of the world-view,
it would be a true saving to abolish the Ministry of Ecclesiastical Affairs and Education;
for what, at bottom, is the purpose of this Ministry
% --- p. 2 ---
\setcounter{footnote}{0}%
\opage{2}but to take care that a certain general world-view,
one that secures society and is therefore
highly necessary for all citizens, small and great,
can be maintained, developed, and brought, in an impressive
manner, to the consciousness of the whole people. But
for the saving to be introduced without danger to
society, a new and great discovery must first be made.
There must absolutely be invented an arcanum, a magic potion,
by whose application one is suddenly able
to implant in all souls, in the mind and heart of all human beings,
an equally clear and living sense of justice, truth,
and fairness, a common unfailing instinct, which tells
everyone that no one can have rights without having
duties, or civil freedom without submitting to the laws,
a warm and strong and heartfelt feeling that we are all
brothers and sisters and are to love one another in all sincerity,
so that no one looks to his own without at the same time
looking to the common good: yes, if anyone could merely
invent such an arcanum, the generations would rejoice
and the spirits of the age would hail the new Messiah, for then the
hardest knots of society would be loosed. What savings of
money, of worry, and of spiritual labor! Then one could
indeed abolish both the Ministry of Ecclesiastical Affairs and compulsory schooling
and the police and the clergy's salaries and the Royal
Theater and the army law; one could in truth introduce
thoroughgoing savings. And what progress in
freedom! If someone wanted to learn something, well and good; if he wanted to
learn nothing, that too was well and good; one could now rely on
his instinct and feeling. If someone wanted, for the sake of salvation,
to have a religion, be it Brahminic, Mohammedan,
Jewish, or Christian, provided he did not want to force
% --- p. 3 ---
\setcounter{footnote}{0}%
\opage{3}his faith on others, well and good; if he wanted no religion
at all, but regarded salvation as a luxury article, that
too was well and good; one could surely rely confidently on his instinct
and feeling. Such perfect freedom
with regard to the religious world-view would of course
bring with it a like freedom with regard to the
moral one. Just as everyone ought to be allowed to dress
according to his taste, so long as he does not go entirely naked, so
everyone should also be allowed to live privately
in his own way; instinct and feeling would surely vouch
that it took place in good understanding with one's fellow men.
If someone wanted to be a Mormon and have seventeen wives, well and
good; if he wanted only a single one, well and good; if
he wanted none, but made do with uncertain income, that
too was well and good. Then freedom would thrive and ``self-government''
would flourish, --- if only we had the magic potion!

Since such a magic potion has not yet, so far as is known, been
invented, the state is compelled to make use of those means of education
which, in accordance with the population's
traditions, faith, and customs, under the existing constitution
prove to be the most effective, and it is undeniably
those that are found in the church and the school. It ought to be expressly
emphasized that the state, although from its point of view it chiefly
must regard these institutions as means
to the maintenance of civil moral order, yet by no means on that account
fails to recognize that they have value and validity
as ends in themselves; but it becomes necessary to
make a sharp distinction between what concerns
the church as church, the school as school, and the state as state,
so that it may be seen from this how much is a matter for
compulsion and what, by the nature of the case, belongs to freedom.

% --- p. 4 ---
\setcounter{footnote}{0}%
\opage{4}For the church as congregation, religion is the sole and
highest end, and faith is free, the proclamation of the Word free;
all congregations are born free congregations, there is absolutely
no exception. None the less the state must regard
religion as a means of education and appraise the
religious world-view from a humane, civic
point of view. For the sake of education it must take care
that the teacher of religion is secured so free and
independent a position over against his hearers that the word of truth
can be heard and the proclaimer of the Word not sink down
to become the client of the majority he is appointed to educate.
The teacher of religion must --- for the sake of education ---
give the state an account, and his work is to be judged, not
by the particular creed to which he adheres,
but by the ability with which he is able to awaken
in souls the living knowledge of truth, the deep reverence
for right and duty that is indispensable for good citizens,
and to make it practical in such a way that
doctrine becomes conviction, and conviction becomes
flesh and blood in the members of the congregation. In
this respect the state ought to exercise a strictly consistent supervision
over all teachers of religion in all congregations of all creeds:
there are --- for the sake of education ---
no exceptions.

For the school, instruction is the main thing. Whether
the emphasis is placed primarily on awakening or on the imparting of knowledge,
the school as school has its own end and
in its own end its freedom. All schools, the schools of the common people
just as much as the learned ones, are in this respect
to be regarded as free schools; no outside
authority ought to interfere arbitrarily in the school's plan: the
% --- p. 5 ---
\setcounter{footnote}{0}%
\opage{5}freer it can go, the better. But the schools are at the same time
institutions of education; however the instruction is conducted,
unless it is quite without character, there will naturally
form a conception of the common human, a
view of life, a generation's world-view. On this
point the state is attentive. What in particular
is the subject of instruction, how much or how little is
read, is not the main question; the question decisive for the state
is how education fares.
If a sound and true and full cognition
of what a capable citizen owes to himself and
to his fatherland is imparted to the young, then all is well; if it does not happen,
then all is bad, even if all the pupils left school
stuffed with learning. The teachers and administrators of the school must on this
point give the state an account; to appeal to freedom
and to give assurances about instinct and feeling
does not help: the state cannot content itself with the uncertain; it
must have full certainty, it must --- for the sake of education ---
exercise an effective supervision over all schools
without exception.

Should it now seem as though we conceded to the state
a guardianship that would have to be oppressive to freedom of spirit
in church and school, we ask it be observed that
at which the guardianship aims is, for the sake of freedom,
limited to the most necessary, the most general,
to something that can be reached by very different
ways. All must admit that a religion which,
however much is preached, cannot even awaken
the civic conscience, is a poor religion, and
a school which, despite lessons, examinations, and tests,
is not able to impress and establish in the
% --- p. 6 ---
\setcounter{footnote}{0}%
\opage{6}minds of the young the first necessary principles of right and
duty, a poor school. In former days the state wanted
too much; it wanted to make schoolboys into dogmatists
and subjects into Christians; that failed. In our
own days, when the state has begun to limit its demands
to the necessary, there is hope that it
will attain precisely what it aims at.

The world-view that the state at any price
wants maintained can be expressed in a few words: it is man's
vocation to live in a society ordered by law,
to obey the laws and to fill his place by a conscientious
discharge of his civic duties; in its generality it is
so colorless that it does not look like much on
paper, but in reality so important that all enthusiasm
for nationality, freedom, and fatherland, without it, becomes
air bubbles and blue vapors. The social world-view
has several enemies that it must combat. It cannot, for example,
coexist with a papal hierarchy, for the hierarchy sets
itself above the civil laws; it cannot coexist
with an ascetic orientation of life that in poverty and renunciation
despises the temporal; it cannot coexist with a selfish
brutality that defies the laws and storms the barriers
of justice in order to take the law into its own hands; it cannot
coexist with a narrow-sighted, low-minded party spirit that only
wants to satisfy the egoism of a single estate, be it a perfume egoism
or a train-oil egoism; and finally it cannot coexist
with the bottomless versatility which, under the banner of
\emph{finer cultivation}, lives by sniffing at everything, chattering about
everything, being done with everything, and doubting everything, except
the intelligent divinity of enjoyment, of diversion, and of opinion.
When these opponents are set aside, the
% --- p. 7 ---
\setcounter{footnote}{0}%
\opage{7}official, civic world-view can be reconciled
with any personal moral conviction whatever.

The personal world-view, which is to bear the
civic one, receives its light and warmth chiefly from the
two focal points, the Deity and the soul. Is there a
God, and what is the soul? This twin question presses
itself upon everyone who begins to think about existence.
Can such questions be answered? Again a
question. Which way is one actually to go in order to
get them answered? Again a question. Then
doubt awakens at the many questions. Our time is a time of doubts.
To doubt everything is, people think, a sign of
superior understanding, it is the beginning of wisdom. And
yet the age, with its busy criticism and its free
thought, has but little joy of all its doubting; it is dispirited,
for doubt consumes; it feels itself faint, disheartened, and suffers
from spiritual hunger. To doubt in order to test properly can be healthy
in a practical as well as a theoretical respect, when
what is to be tested is of such a nature that the testing can
end the doubt. But to doubt personally about one's own self,
about the meaning of life and the origin of the world, to nurse doubt,
to brood over doubt, to doze in doubt, is unhealthy; such
doubting is sickness. The present suffers from the doubting sickness; the
pulse of thoughts is restless, the world has cramps, the race has
fever; it is the worms of self-contradiction that gnaw at the roots
of the heart. Under violent chills the most absurd
conceptions find a voice, and one contradiction chases the
other. Modern man literally does not know
what he is to take himself for, an animal, a piece of God,
or both. On the one hand a language is used
in criticism and philosophy as if human knowledge had
% --- p. 8 ---
\setcounter{footnote}{0}%
\opage{8}no innate, insurmountable limit; the riddles of the world
are solved, the old night of possibilities must out with its secrets:
existence has no mystery, the Absolute
itself lies open, transparent to the thinker's
gaze. One can well understand that, if it is really
true that spirit first becomes spirit in man; one
can understand it when it is held consistently that the Deity,
this principle of all-nature with its blind wisdom,
whose light is darkness, is a principle of consciousness that
strives upward toward self-consciousness, but cannot become conscious of
itself until, with the help of matter, it has built
itself a human brain. From this point of view man
is then the self-conscious Deity, and the self-conscious
Deity must in science --- gradually, of course,
through many exertions and only by approximation! ---
be able to see through its own unconscious work. Now
one would suppose that, in the name of knowledge, consciousness was meant in earnest,
that philosophical self-consciousness, the
absolute form of the absolute content, must be the in itself
highest, most glorious, most perfect of all; but no! it
is not so meant. Man is, this must above all
not be forgotten, an animal, a scientific animal. Only as
animal has he consciousness, only by virtue of consciousness has
he self-consciousness: the animal's self-consciousness is as perishable
as the animal itself. One sees how the illumination
shifts. \emph{First the conscious was to be higher than
the unconscious}; the Absolute had to go the long
way around before it brought forth consciousness, for it sought
knowledge, and only in consciousness was there knowledge. But
then it suddenly dawns on us: consciousness has
only a limited knowledge, dependent on vibrations of the brain;
% --- p. 9 ---
\setcounter{footnote}{0}%
\opage{9}if disorder comes into the vibrations of the brain --- which unfortunately
can easily happen --- then disorder comes into
knowledge too; if a thinker got water on the brain, the knowing
Absolute would have to get water on the brain too; the
eternally divine knowledge cannot be served by that. What
is knowledge to do, then? Emancipate itself from all heads,
all brains and vibrations of the brain, that is, from consciousness!
that was a clever turn. \emph{Now the unconscious is
higher than the conscious}. The Absolute then has no
trouble with vibrations of the brain; it can do without brains,
for it can do without all consciousness. Were all thinkers
to land in the madhouse, the Absolute would not land in the
madhouse, for it can do without all thinkers. Perfect
knowledge is a knowledge in the unconscious, a blessed
knowledge. How comforting, not merely for the Absolute, but
also for the animal with the brain. Now the knowing animal can
die when it will; it is true that the brain goes with it,
for the brain is consciousness; but what knows in the brain
suffers no harm, on the contrary! the brainless knowledge
is a blessed knowledge, blessed in the unconscious. That here are two
irreconcilable world-views no one will be able to deny, and
yet they must be reconciled and melted together into one, for
otherwise dualism comes.

It takes artistry to draw
a square circle; and it takes acuteness and
a fine knowledge of the history of philosophy to
reconcile two views as irreconcilable as the above.

If one starts from the assumption that brain and
consciousness in man must stand and fall together,
then one has only a choice between one of two things. Either
to admit that there is a divine consciousness,
% --- p. 10 ---
\setcounter{footnote}{0}%
\opage{10}infinitely exalted above the human, since
there is a divine knowledge infinitely exalted
above the human; in this case the limit of human knowledge, conditioned by the brain,
is recognized once and for all, and there is a mystery of existence.
Or openly to declare that in the essence of things there are no
thoughts at all, since in the things themselves there is no consciousness,
and that unconscious thoughts are impossible. --- There
are obscure thoughts in the soul, because there is an obscure
consciousness and an obscure thinking. The thought is obscure
when the consciousness is obscure; but no thought, be it
clear or obscure, is possible without a corresponding consciousness.
In order to think the unconscious we must as
necessarily use our consciousness as in order to think the
conscious. We say of a stone that it is unconscious, but
we say it by virtue of our consciousness. He who has never
seen the light has no conception of darkness, and
he who has never become conscious of anything does not know
that there is anything unconscious, or what the unconscious is:
he has no knowledge at all.

If some day a prince or a rich Englishman took
it into his head to put up a couple of million pounds sterling as
a prize for the answer to the question: \emph{How does
human consciousness go about
thinking a perfect knowledge without all consciousness}?
then it would not be a thinker of ordinary organization,
it would have to be a genuine Absolute, who won
the prize. It is an easy matter to subject the forces of nature
to an unconscious will and to speak of nature's unconscious
purpose, when by that nothing is to be thought of other than a
figurative transference of what stems from consciousness
to that which has no consciousness. If one calls without
% --- p. 11 ---
\setcounter{footnote}{0}%
\opage{11}further all striving will, then there is will everywhere, will
in the stone to fall, will in the moon to approach the
earth, and so on. For that matter one might as well
say that in the locomotive of a railway there is an unconscious
will to pull the carriages, joined with an unconscious
representation of the next station to be reached.
But these are only figures of speech. Precisely the fact
that nature works as if according to purpose, although it has
no purpose, that it so wonderfully reminds one of an All-knowing One,
though it has no knowledge at all, is the great
riddle. When human knowledge, in order to be able to
contain the Absolute, absolutely must be rid of consciousness
and yet cannot get rid of consciousness, then it is badly
off; doubts rise up, contradiction torments, and the anguish
can become Tantalean.

Everyone who, for the sake of personal conviction,
feels the need to form a world-view for himself had better
see to it that he has not only his thoughts free, but also his mind
free. Without free thoughts he will on many points
cling to what has hitherto been accepted, to deep-rooted presuppositions
that are not tested, and since he nevertheless adds a good deal
of his own, the working through of the whole
will be so imperfect that coherence is lacking.
But even if the work of thinking is carried out with such
consistency that an exacting criticism will hardly
be able to point out essential disagreements between the individual
parts of the edifice, that is still no full proof
of its truth as a world-view. A world-view
can, in formal respect, have scientific
consistency and good coherence and in reality yet
be false, thoroughly false, indeed absurd to the point of madness,
% --- p. 12 ---
\setcounter{footnote}{0}%
\opage{12}when the thinker's mind is ensnared and unfree. We can cite
a magnificent example of this from the most recent philosophy in
Germany.

``The Philosophy of the Unconscious'';\footnote[1]{Philosophie des Unbewussten. Versuch einer Weltanschauung von E. v. Hartmann, Dr. phil. Berlin 1869.} a book that has already, and
with reason, aroused considerable attention in
German literature, is more than a well-written book; it
is a work of the mind, carried through in fine strokes and with acute analysis of the
most difficult problems of natural science,
especially of physiology, that lies before us here. Upon
an appraisal of the work's scientific design and peculiar
point of departure we do not wish in the present context
to enter; only from the world-view expressed
in it do we cite a few features as a sample. --- From
Schelling's first youthful appearance with his
intellectual intuition of the Absolute until the latest
times, speculative philosophy, in equivocal accord
with natural science, has labored to clarify for itself the absolute
being. It has meanwhile been difficult for the intelligent, co-philosophizing
public to follow and to keep up with
how it actually stood with this Absolute.
At one time it was presented as a logical all-knowing
Deity: God in Logic, before he created the world,
at another as an undisclosed, obscure principle, the indwelling
ground and cause of all existence. Now at last
the equivocation is removed; the visage of the goddess of truth is seen completely
unveiled. The Philosophy of the Unconscious, which does not one-sidedly
keep to speculation, but accurately revises the facts of
experience and takes a vigorous grip on
% --- p. 13 ---
\setcounter{footnote}{0}%
\opage{13}materialism, has fathomed the secret. The great
specter that has so long haunted speculative immanent
brains has been conjured forth, the Absolute shows itself in
broad daylight and the Absolute is? --- the pure Nothing, the
eternally blessed ``gar Nichts.''

The riddle of the world is that an unconscious willing from the beginning
has driven an equally unconscious representing from
non-being into being, for, as it says on p. 607, ``in itself
the representation had no drive and no interest in
passing from non-being over into being, consequently before the onset
of willing no representing was actual, consequently there was
before the coming-to-be of the world neither willing nor representing
i.e.\ ``gar Nichts''. As long as the will continues, so long
will the process and its appearance in consciousness,
the world, continue; when therefore there is once to be
no world any more, then there may no more be
any willing or any representing (since the unconscious
representation becomes actual only insofar as the interest
of the will requires it) i.e.\ there will again be ``gar
\emph{Nichts}.'' With this the following description
of the Deity (p. 606) also agrees: ``The All-One Unconscious
is all-knowing and all-wise, and consequently cannot become wiser
than it already is; it has, as Aristotle also says,
no memory, and so it can learn nothing from the experiences it
otherwise made in the world. Consequently, when
the world has once ceased to be, it has remained exactly
the same as it was before the creation of the world; as blessed
as it was before, so it is now again, neither more nor
less; the world-process can in no way help
it to a greater blessedness than it had before\dots{}.''

If the Deity is blessed in Nothing, it follows of
% --- p. 14 ---
\setcounter{footnote}{0}%
\opage{14}itself that souls too must become blessed in Nothing. The hope
of an individual continuance of the soul after death is declared
outright to be an illusion. ``Thereby the main nerve of the Christian
promises is severed, for man is at bottom
concerned only with his own dear ego;
what good to me is the greatest future blessedness, if
I cannot perceive and enjoy it!''

One gathers from this what the world is, what life is
worth, and what all human poetizing and striving has to
mean; it is disappointment, conceit, illusion, all of it.
``The Philosophy of the Unconscious'' sets out for us three stages of consciousness
corresponding to three stages on the way of illusions. The first
stage of illusion is marked by the individual's imagining
that, in virtue of a given state of the world, he can attain happiness
in the earthly life. That it is a deception follows from the fact
that all pleasure is only a condition for displeasure, all enjoyment for
suffering, all joy for sorrow and pain. If human beings
would only make it clear to themselves that joy stands to
the sorrow undergone as the cessation of pain to the torments
of the rack, they would soon be cured of the
spreading sickness for enjoyment. ``What one nowadays calls
the specter of pauperism is nothing other than the consciousness
emerging in the masses that the struggle with want,
sorrows and their alleviation lies wholly on the negative
(pain-) side of the zero point of sensation, whereas
formerly, when the poverty of the masses was ten times greater, this
consciousness was lacking, and the people bore their poverty as
by the grace of God!'' On the second stage of illusion the individual
imagines that he is to attain happiness in a life beyond,
after death. That is the Christian deception. ``Jesus's world-view
was too naive and childlike to separate soul
% --- p. 15 ---
\setcounter{footnote}{0}%
\opage{15}and body,'' which is why ``the resurrection of the flesh'' is
included in the third article of the creed. That it is an illusion goes
without saying; indeed, there are even passages in Scripture that
intimate that all consciousness, like all time, ceases in eternity.
The third and last stage of illusion is that on which
we are riding as on a railway. The conceit is that
happiness lies in a great political and social future. One
expects salvation from an improved order of society, from
progress in industry and science: a monstrous deception!
The end is that the old, worn-out human race is
tired unto death, tired of the labors and of the enjoyments,
without heirs and without love, tired of everything. Then
there will come that exalted melancholy which is usually
found in geniuses and in spiritually eminent old men; then
it will come, the pure clearness of heaven, that cool divine
calm that wafts through Spinoza's Ethics, where the passions
are swallowed in the abyss of reason, because they are clearly and
distinctly grasped in ideas. Then illusion is dead, hope
burnt out, for what should one still hope for? Then
the dead-tired humanity recognizes only one thing that is
blessed, rest from sorrow and want, painlessness in the pure
Nothing.

Is this now a comfortless, desperate world-view?
By no means! that is to say, when the mind is free and unbiased like
thought. Toward an objective scientific theory the
mind must be free. That the mind is free can mean that the thinker
does not make the theory his personal conviction, but
has an interest solely in examining what it is built on,
and how far it hangs together in all its parts. It
can also mean that he makes the theory his personal
world-view, provided he only knows within himself that
% --- p. 16 ---
\setcounter{footnote}{0}%
\opage{16}he does it by virtue of a choice. In this there is nothing comfortless.
For why does he \emph{choose} the view, unless
because he feels that it is precisely a view for him, a
world-view that deeply and essentially satisfies his
inner need, a view in which his weary soul finds
rest. If someone, on the other hand, though in his inmost being
he feels a need for a living truth, a true cognition
of eternal life, lets himself be overpowered by such a wisdom of death
and annihilation, because he is impressed
by the supposedly objective ``Wahrheit,'' then he is in
truth comfortless. But the comfortless does not lie in
the truth; no, it lies in his own self-deception. His
heart is sick, his mind is ensnared; he imagines
that he stands face to face with the most terrible of all
powers, the truth itself, and then it is not the truth at all!
The theory's objective specter is precisely the illusion,
indeed the illusion of illusions, the most powerless being of all,
the shadow of a Nothing: the objective ``Wahrheit,'' separated from divine and
human consciousness, is --- ``gar
Nichts.''

To form a world-view for oneself is by no means the
same as to form a correct representation, a concept as
certain and distinct as possible, of the outward
real world. The cognition of the laws of things, the
course of history, and the coherence of the world can in an observer
be highly developed, and the presuppositions on which
a personal conviction is to be built nevertheless be
extremely obscure, deficient, confused. In order to cognize
the true, we must cognize the real and learn to
regard the facts, be they of nature or of history,
in their own light; that is certain enough. A world-view
so shy of light that it dares not look reality
% --- p. 17 ---
\setcounter{footnote}{0}%
\opage{17}straight in the eye, so pampered a subjectivity that it slinks
out of the way of science because it is afraid to meet
the objective, so selfishly conceited that it positively bears
a grudge against natural science, philosophy, and criticism, because it has
a suspicion that these airers might easily take it into their heads
to make use of the broom and tidy up a little in its
nursery, such a world-view is naturally
without value for those who seek the light and love the truth.
But the cognition of reality is only the introduction to the matter,
not the matter. To a personal world-view there belongs
necessarily a personal self. He who forgets his own
self forgets what is most important of all; bewitched by the objective,
he has already at the point of departure begun by
deceiving himself and ensnaring himself in illusion.

Now it is doubtless so that a world-view cannot
be arbitrary, that the human self cannot
create itself as it will, because it cannot remake the
real world and alter the order of nature; there is a
fundamental law of life that holds for all that lives without exception;
only one revelation of the personal self can,
in an eternal sense, be the perfectly true. But it
is also certain that each single human being's cognition of
the self is conditioned by his self-knowledge, and this self-knowledge
by his self-development. Since it is now notorious
that individuals in society find themselves at very
different stages of self-development and self-knowledge, and
since it is impossible that a personally convincing world-view
can gain entrance with anyone without going
the way of self-knowledge: it follows from this that the many
mutually conflicting world-views, materialism, nihilism,
atheism, pantheism, rationalism, Judaism,
% --- p. 18 ---
\setcounter{footnote}{0}%
\opage{18}Christianity, and so on, must all have full right to exist
and to express themselves in a society ordered on the basis of freedom.

What then has the state to do in relation to the many personal
world-views? Give up an old
prejudice; renounce a thousand-year-old, deep-rooted illusion, and
renounce it with the same sincerity with which a believing
candidate for baptism renounces the Devil and all his works and all
his deeds, --- the old soul-corrupting illusion
that a uniform, fully personal world-view should
be something that could be forced by objective state laws upon
all or at least upon most; that there should be an objective
religious science, a science of faith, that could come
to the aid of the state authority in getting the objective
truth generally introduced; that natural science, philosophy, and
criticism should be regarded only as provisionally tolerated, ``heathenishly
puffed-up sciences,'' which can only obtain the forgiveness of
sins when they again condescend to
kneel before theology and let themselves be baptized upon dogmatics.
By thoroughly renouncing this old illusion and vigorously
taking care of the morally civic education of the young,
by supervision over the whole of religious as well as humane
instruction, the state will attain its end and feel secure.

The many world-views are personal powers of life;
they fight with the weapons of life for the right of life.
As the struggle goes, it will show which among
them has the strength, because it has the conditions.
It shall show that they are not equally strong, because they
are not equally true; it shall show in facts
that there is only one that is eternally strong, because
there is only one that is eternally true, a single one, that
pearl of the Gospel, which must be bought at the dearest price.
\signature{R. \emph{Nielsen}.}

\end{document}
